% s-075-p075-m002; original PDF 75
\Aphorism{s-075-p075-m002}{65}
Der Egoist hat schon drei Viertel seines inneren Menschen gegen sich; um so eher darf er seine äusseren Widersacher geringschätzen. Wenn es ursprünglich vielleicht Mühe gekostet hat, ein paar sociale Instincte in das Individuum hineinzuzüchten, so kostet es heute umgekehrt Mühe, nicht social, sondern individualistisch zu empfinden. Es ist natürlich hundertmal leichter, zu entsagen und sich mit dem Augenblicksvortheil des guten Gewissens zufrieden zu geben, als seiner ganzen Mitwelt und seiner eigenen Vergangenheit zum Trotz den Weg zum eigenen Ich, einen ungebahnten Weg zu einem ungewissen Ziele, zu gehen. Niemandem sei es verdacht, wenn er von diesem Weg bei Seite schleicht und seine kostbare Verantwortlichkeit lieber in altruistischen Werthen anlegt; nur mögen diese Abtrünnigen sich nicht einreden, sie seien die Opfernden und der Egoist sei der Schlaue und Bequeme.
\par

